% !PRINTED 60
\weberpage{60}

\begin{center}
\textbf{§ 17. Die Funktionen $T$.}
\end{center}

Wenn doppeltperiodische Funktionen $\varphi(u)$ von der $m$ten Ordnung existieren, so müssen es gebrochene Funktionen sein. Setzen wir sie also nach \S\,15, 10. in der Form

\weq{1}{\varphi(u) = \frac{G_1(u)}{G(u)},}

wenn $G_1(u)$ und $G(u)$ ganze Funktionen ohne gemeinsamen Nullpunkt sind, so ist

\[
\frac{G_1(u+\omega_1)}{G(u+\omega_1)} = \frac{G_1(u+\omega_2)}{G(u+\omega_2)} = \frac{G_1(u)}{G(u)},
\]

und darin haben die drei Funktionen $G(u)$, $G(u+\omega_1)$, $G(u+\omega_2)$ dieselben $m$ Nullpunkte, nämlich die Unstetigkeitspunkte von $\varphi(u)$. Ihre Quotienten sind also Einheitsfunktionen, und nach \S\,15, 9. ist also, wenn $g_1(u)$, $g_2(u)$ \gesperrt{ganze Funktionen} sind:

\weq{2}{\begin{aligned}
G(u+\omega_1) &= e^{g_1(u)}\,G(u),\\
G(u+\omega_2) &= e^{g_2(u)}\,G(u),
\end{aligned}}

und die Funktion $G_1(u)$ genügt den nämlichen Gleichungen.

Ist

\weq{3}{\varphi(u) = \frac{T_1(u)}{T(u)}}

eine zweite Darstellung von $\varphi(u)$ von der Form (1), so ist

\weq{4}{T(u) = e^{g(u)}\,G(u),}

worin $g(u)$ wieder eine ganze Funktion ist, und wenn also nach (2)

\weq{5}{\begin{aligned}
T(u+\omega_1) &= e^{l_1(u)}\,T(u),\\
T(u+\omega_2) &= e^{l_2(u)}\,T(u)
\end{aligned}}

ist, so ist

\[
\begin{aligned}
g_1(u) &= l_1(u) + g(u) - g(u+\omega_1),\\
g_2(u) &= l_2(u) + g(u) - g(u+\omega_2).
\end{aligned}
\]

Es kommt also alles darauf an, daß wir unter irgend einer passenden Annahme über die ganzen Funktionen $l_1(u)$, $l_2(u)$ ganze Funktionen $T(u)$ bilden können, die den Bedingungen (5) genügen und in $m$ beliebig gegebenen Punkten des Periodenparallelogramms verschwinden. Durch die Quotienten zweier solcher Funktionen mit denselben $l_1(u)$, $l_2(u)$ können wir dann alle doppeltperiodischen Funktionen $\varphi(u)$ darstellen, und die allgemeineren Funktionen $G(u)$, die dasselbe leisten, erhält man nach (4) durch Hinzufügung eines \gesperrt{willkürlichen Einheitsfaktors}.
